\scanpage{93}\mspage{T-7}

5. \uline{Forme canonique} ( $p=2$)

On a ici $ua'_{1}=a_{1}$ (cf. p.T-3). On distingue deux cas :

\uline{Cas général} : $a_{1}\ne 0$ (equivaut à $j\ne 0$).

En choisissant convenablement $u,r,t$, on s'arrange pour que $a_{1}=1$, $a_{3}=0$, $a_{4}=0$. D'où la forme canonique :
\[
 y^{2}+xy=x^{3}+a_{2}x^{2}+a_{6}\qquad\qquad \pi=dx/x\ .
\]

On a :
\[
 b_{2}=1,\ b_{4}=b_{6}=0,\ b_{8}=a_{6},\ c_{4}=1,\ \handwritten{\Delta}=a_{6}\ ,\
 j=\frac{1}{a_{6}}\ .
\]

Isomorphismes :
\[
 x\circ f=x'\ ,\quad y\circ f=y'+sx'\ ,
\]
avec :
\[
\begin{aligned}
 a'_{2}&=a_{2}+s^{2}-s\\
 a'_{6}&=a_{6}\ .
\end{aligned}
\]

Deux courbes $C$ et $C'$ de même $j$ sont isomorphes si et seulement si $a'_{2}-a_{2}$ est de la forme $s^{2}-s$ ; cette condition est vérifiée sur une extension séparable de $k$ de degré $\handwritten{\leqslant}2$.

Le groupe d'automorphismes de $C$ a deux éléments, correspondant à $s=0,1$.

Courbe type :
\[
 y^{2}+xy=x^{3}+a_{6}\ .
\]

\uline{Cas}\quad $a_{1}=0$\quad (équivaut à $j=0$).

En choisissant convenablement $r$, on s'arrange pour que $a_{2}=0$. D'où la forme canonique :
\[
 y^{2}+a_{3}y=x^{3}+a_{4}x+a_{6}\qquad\qquad \pi=dx/a_{3}\ .
\]

On a :
\[
 b_{2}=b_{4}=0,\ b_{6}=a_{3}^{2},\ b_{8}=a_{4}^{2}\ ,\
 \handwritten{\Delta}=a_{3}^{4}\ ,\ j=0.
\]
% archive sheet number: 92

